\scanpage{350}
% @ head of sheet
\mspage{74}
% @ top, lines 1–2
toute fibre de $p$ est $\inserted{\unsure{\sim|\simeq}}\underrightarrow{\lim}_{I}$ des fibres correspondantes des $p_\alpha$. De cette façon, on trouve

% @ upper third, lines 3–6
\uline{Proposition} 5.1. Supposons que $W$ satisfasse l'axiome \inserted{des $W$-fibrations} $\mathrm{Loc}(7\,D_{W_{\unsure{\omega|W|w}}})$ (ce qui est donc le cas s'il satisfait un des axiomes \inserted{plus forts} $\mathrm{Loc}(7\,a\ \text{ou}\ b\ \text{ou}\ c)$).

% @ middle, lines 7–12
Pour que $W$ satisfasse à $\mathrm{Loc}(8)$ (resp. à $\mathrm{Loc}(8\mathrm{bis})$), il f. et s. \deleted{qu'il} \deleted{soit} que l'ensemble des objets de $\mathrm{Cat}$ qui sont $W$-asphériques soit stable par $\underrightarrow{\lim}$ filtrantes (resp. par $\underrightarrow{\lim}$ \deleted{\unsure{f} \unsure{filtrantes}} de suites \unsure{dénombrables}).

% @ lower third, lines 13–16
\deleted{\unsure{Corollaire}} \inserted{5.2.} Notons que la catégorie $\mathrm{Cat}$ est un \unsure{paratopos|prétopos}, - a fortiori elle est \uline{accessible}, \inserted{cf $\overline{\underline{\mathrm{XVIII}}}$ \S\,8),} et il en est donc de $=$ de $\underline{\mathrm{Fl}}(\mathrm{Cat})$. Ainsi on a dans $\mathrm{Fl}(\mathrm{Cat})=\mathrm{Ob}\,\underline{\deleted{\unsure{\mathrm{E}}}\mathrm{Fl}}(\mathrm{Cat})$ la notion
% @ foot, lines 17–21
de \uline{partie accessible}. On dira donc qu'une partie $W$ de $\mathrm{Fl}\,\mathrm{Cat}$ (plus gén. d'une catégorie $\mathrm{Fl}(M)$, où $M$ est une catégorie accessible) est \uline{accessible}, si elle l'est \deleted{dans la} \unsure{p.r.} à la catégorie ambiante
% archive sheet number: 349
